\noindent\textbf{Elección del procedimiento.} El equilibrio geostrófico iguala gradiente horizontal de presión y Coriolis. En módulo, $\rho fV_g=\Delta p/\Delta n$. Se calcula primero $f=2\Omega\sin\varphi$, después $V_g$ y por último el Rossby $V_g/(fL)$ con la escala horizontal elegida.

\begin{enumerate}

\item Usar latitud en radianes y la componente vertical de la rotación.

\[\fitmath{\mathrm{fcalc}=2\cdot \Omega\cdot \sin\left(\varphi\right)}\]

\[\fitmath{\mathrm{fcalc}=2\cdot (7{,}2921\times10^{-5})\cdot \sin\left((0{,}7853981634)\right)\approx 0{,}0001031259\,\mathrm{1/s}}\]

\item Convertir 4 hPa a 400 Pa y 400 km a 400000 m.

\[\fitmath{\mathrm{gradp}=\frac{\mathrm{dp}}{\mathrm{dn}}}\]

\[\fitmath{\mathrm{gradp}=\frac{(400)}{(400000)}\approx 0{,}001\,\mathrm{Pa/m}}\]

\item Despejar la velocidad del balance geostrófico.

\[\fitmath{\mathrm{Vg}=\frac{\mathrm{gradp}}{\rho\cdot \mathrm{fcalc}}}\]

\[\fitmath{\mathrm{Vg}=\frac{(0{,}001)}{(1{,}2)\cdot (0{,}0001031258672)}\approx 8{,}08074\,\mathrm{m/s}}\]

\item La escala de 1000 km del enunciado permite comparar inercia con Coriolis.

\[\fitmath{\mathrm{Ro}=\frac{\mathrm{Vg}}{\mathrm{fcalc}\cdot \mathrm{Lscale}}}\]

\[\fitmath{\mathrm{Ro}=\frac{(8{,}08074013)}{(0{,}0001031258672)\cdot (1000000)}\approx 0{,}07835803\,}\]

\end{enumerate}
